\chapter{Detailed Proofs}
\label{app:proofs}

This appendix gives complete proofs of every numbered mathematical claim in
Chapters~\ref{ch:success}--\ref{ch:censored}, together with precise statements and proofs of results
the main text uses in display form only: the selection identities of Chapter~\ref{ch:us}, the
instability of log-link count feedback behind Warning~\ref{warn:poi-metastable}, and the explicit
$\kappa$--$\theta$ ridge computation behind Proposition~\ref{prop:cen-ridge}. Each claim is restated
compactly --- keeping its original number, which the restatement references --- and then proved in
full: every interchange of limit, integral, and expectation is justified, and every appeal to the
literature is explicit. Where the fully general theorem is genuinely beyond the scope of this manual
(the general marked-space Watanabe theory, stability of nonlinear Hawkes processes, the do-calculus
in full), we prove the special case the manual actually uses and cite the source for the general
statement. Helper results proved here carry appendix numbers and labels prefixed \texttt{app}.

\section{Standing conventions and two lemmas}
\label{app:proofs-standing}

Throughout, $(\Omega, \mathcal{F}, \Prob)$ carries a filtration $(\Hist_t)_{t \ge 0}$ satisfying the
usual conditions, and $N$ is a \emph{simple} point process adapted to it: a right-continuous
counting path with unit jumps at strictly increasing event times $t_1 < t_2 < \cdots$ having no
finite accumulation point. We assume $N$ admits a nonnegative predictable intensity $\lambda$ in the
sense of Definition~\ref{def:poi-intensity}, with
\begin{equation}
\label{eq:app-integrable}
\E[\Lambda(T)] \;=\; \E\Big[\int_0^T \lambda(u)\,\dd u\Big] \;<\; \infty
\qquad \text{for every finite } T,
\end{equation}
and write $M = N - \Lambda$ for the compensated process.

\begin{lemma}[Compensator calculus]
\label{lem:app-mart}
Under \eqref{eq:app-integrable}: (i) $M$ is a martingale on every $[0,T]$; (ii) for every bounded
predictable real- or complex-valued process $H$, the integral $t \mapsto \int_0^{t} H(s)\,\dd M(s)$
is a martingale on $[0,T]$; in particular it has mean zero.
\end{lemma}

Lemma~\ref{lem:app-mart} is the defining property of the compensator together with the elementary
theory of Stieltjes integration against point-process martingales; we use it without proof and refer
to \cite{bremaud1996} (the integration theory recalled there) and \cite[Ch.~7]{daley2003}. Complex
integrands are covered by applying the real statement to real and imaginary parts. The special case
we use most often we prove in full.

\begin{lemma}[Unbiasedness against deterministic integrands]
\label{lem:app-fubini}
Let $f : [0,T] \to [0,\infty]$ be Borel measurable. Then $\E \int_0^T f(s)\, \dd N(s) = \E \int_0^T
f(s)\, \lambda(s)\, \dd s$, both sides possibly infinite. For bounded signed measurable $f$ the same
identity holds with both sides finite.
\end{lemma}

\begin{proof}
Define two measures on the Borel subsets of $[0,T]$ by $\mu_1(B) = \E \int_0^T \ind{s \in B}\, \dd
N(s)$ and $\mu_2(B) = \E \int_0^T \ind{s \in B}\, \lambda(s)\, \dd s$; both are finite by
\eqref{eq:app-integrable}. For $B = (a,b]$, Lemma~\ref{lem:app-mart}(i) gives $\E[N(b) - N(a)] =
\E[\Lambda(b) - \Lambda(a)]$, so $\mu_1 = \mu_2$ on the $\pi$-system of half-open intervals, and
both vanish on $\{0\}$. Two finite measures agreeing on a generating $\pi$-system agree on the
generated $\sigma$-algebra (Dynkin), so $\mu_1 = \mu_2$. The identity extends to simple nonnegative
$f$ by linearity; to measurable $f \ge 0$ by monotone convergence along simple $f_n \uparrow f$,
each side being an integral against $\mu_i$ by Tonelli; and to bounded signed $f$ by splitting $f =
f_+ - f_-$.
\end{proof}

\section{Proofs for Chapter~\ref{ch:success}}
\label{app:proofs-ch3}

\paragraph{Claim (Proposition~\ref{prop:success-arc}).}
\emph{The arc growth rate \eqref{eq:success-arc} is defined for all sign patterns, bounded in
$[-2,2]$, antisymmetric under swapping $F_y \leftrightarrow F_{y+h}$, and for $F_y, F_{y+h} > 0$
satisfies $g = \log(F_{y+h}/F_y) + O(\log^3(F_{y+h}/F_y))$ near zero growth.}

\begin{proof}
Write $a = F_y$, $b = F_{y+h}$, $g = (b-a)\big/\tfrac12(|a| + |b|)$.

\emph{Definedness.} The denominator vanishes iff $|a| + |b| = 0$ iff $a = b = 0$, the case covered
by the convention $g = 0$.

\emph{Bounds, with equality characterization.} By the triangle inequality $|b - a| \le |a| + |b|$,
so $|g| \le 2$. Squaring, $|b-a|^2 \le (|a|+|b|)^2$ is equivalent to $-2ab \le 2|a|\,|b|$, with
equality iff $ab = -|a|\,|b|$, i.e.\ iff $ab \le 0$. Thus $|g| = 2$ exactly for sign-crossing (or
from-zero) transitions, the reading of Example~\ref{ex:success-arc}.

\emph{Antisymmetry.} Swapping $a$ and $b$ negates the numerator and fixes the denominator.

\emph{Expansion.} For $a, b > 0$ set $u = \log(b/a)$, so $b = a e^{u}$ and
\[
g \;=\; \frac{a(e^{u} - 1)}{\tfrac12\, a\,(e^{u} + 1)} \;=\; 2\,\frac{e^{u} - 1}{e^{u} + 1}
\;=\; 2\tanh(u/2).
\]
Since $\tanh$ is analytic with bounded derivatives of every order on $\mathbb{R}$, Taylor's theorem
with Lagrange remainder gives $\tanh x = x - x^3/3 + O(x^5)$ uniformly on compacts, hence $g = u -
u^3/12 + O(u^5)$ and in particular $g - u = O(u^3)$: arc growth agrees with log growth to second
order at zero growth.
\end{proof}

\section{Proofs for Chapter~\ref{ch:xgb}}
\label{app:proofs-ch4}

\paragraph{Claim (Proposition~\ref{prop:xgb-pinball}).}
\emph{Let $Y$ have distribution function $F$ with $\E|Y| < \infty$ and let $\varphi(a) =
\E[\rho_\tau(Y - a)]$ with $\rho_\tau$ the pinball loss \eqref{eq:xgb-pinball}. Then $a^*$ minimizes
$\varphi$ if and only if $F(a^{*-}) \le \tau \le F(a^*)$; if $F$ is continuous and strictly
increasing, the minimizer is unique and equals $F^{-1}(\tau)$.}

\begin{proof}
Since $0 \le \rho_\tau(u) \le \max(\tau, 1-\tau)\,|u|$, $\varphi$ is finite everywhere. Fix $y \in
\mathbb{R}$ and let $\psi_y(a) = \rho_\tau(y - a)$: on $\{a < y\}$ it equals $\tau(y-a)$, on $\{a >
y\}$ it equals $(1-\tau)(a-y)$, so $\psi_y$ is Lipschitz with constant $\max(\tau, 1-\tau)$ and
derivative $\ind{y < s} - \tau$ at every $s \ne y$. By the fundamental theorem of calculus for
Lipschitz functions, for $a < b$, $\psi_y(b) - \psi_y(a) = \int_a^b (\ind{y < s} - \tau)\, \dd s$.
Taking expectations and applying Fubini --- the integrand is bounded by $1$ on the product of a
probability space with $[a,b]$ ---
\begin{equation}
\label{eq:app-pinball-key}
\varphi(b) - \varphi(a) \;=\; \int_a^b \big( \Prob(Y < s) - \tau \big)\, \dd s
\;=\; \int_a^b \big( F(s) - \tau \big)\, \dd s ,
\end{equation}
using $\Prob(Y < s) = F(s^-) = F(s)$ off the (countable, Lebesgue-null) set of discontinuities.

Let $a^*$ satisfy $F(a^{*-}) \le \tau \le F(a^*)$. For $b > a^*$, monotonicity gives $F(s) \ge
F(a^*) \ge \tau$ on $(a^*, b)$, so \eqref{eq:app-pinball-key} yields $\varphi(b) \ge \varphi(a^*)$;
for $b < a^*$, $F(s) \le F(a^{*-}) \le \tau$ on $(b, a^*)$, so $\varphi(a^*) - \varphi(b) \le 0$.
Hence $a^*$ is a global minimizer. Conversely, if $F(a) < \tau$, right-continuity gives $\delta > 0$
with $F < \tau$ on $[a, a+\delta)$, so $\varphi(a + \delta) < \varphi(a)$ and $a$ is not a
minimizer; if $F(a^-) > \tau$, there is $\delta > 0$ with $F > \tau$ on $(a - \delta, a)$, so
$\varphi(a - \delta) < \varphi(a)$. When $F$ is continuous and strictly increasing, the condition
$F(a^-) \le \tau \le F(a)$ reads $F(a) = \tau$, which has the unique solution $a = F^{-1}(\tau)$.
\end{proof}

\begin{lemma}[Quantile of exchangeable scores]
\label{lem:app-rank}
Let $V_1, \dots, V_{n+1}$ be exchangeable real random variables, $\alpha \in (0,1)$, and $k = \lceil
(1-\alpha)(n+1) \rceil \le n$. Let $Q$ be the $k$-th smallest of $V_1, \dots, V_n$. Then
$\Prob(V_{n+1} \le Q) \ge 1 - \alpha$, and if the $V_j$ are almost surely distinct, also
$\Prob(V_{n+1} \le Q) \le 1 - \alpha + 1/(n+1)$.
\end{lemma}

\begin{proof}
For $j \le n+1$ let $R_j = \#\{i \ne j : V_i < V_j\}$, computed in the pooled sample of size $n+1$.
First, $\{V_{n+1} > Q\} \subseteq \{R_{n+1} \ge k\}$: if $V_{n+1}$ exceeds the $k$-th smallest of
the first $n$ values, then at least $k$ of those values are $\le Q < V_{n+1}$. Second,
deterministically $\#\{j : R_j \ge k\} \le n+1-k$: if $R_j \ge k$ then $V_j$ strictly exceeds at
least $k$ pooled values, hence exceeds the $k$-th smallest pooled value $W$; and at least $k$
indices have $V_i \le W$, so at most $n+1-k$ indices can satisfy $V_j > W$. Third, by
exchangeability the events $\{R_j \ge k\}$ have a common probability (swapping coordinates $j$ and
$n+1$ maps one to the other and preserves the joint law), so
\[
\Prob(R_{n+1} \ge k) \;=\; \frac{1}{n+1} \sum_{j=1}^{n+1} \Prob(R_j \ge k)
\;=\; \frac{\E\,\#\{j : R_j \ge k\}}{n+1} \;\le\; \frac{n+1-k}{n+1} \;\le\; \alpha ,
\]
the last step because $k \ge (1-\alpha)(n+1)$. This proves the lower bound. For the upper bound with
almost surely distinct values: $\{V_{n+1} \le Q\}$ then equals, up to a null set, $\{V_{n+1} < Q\}
\subseteq \{R_{n+1} \le k-1\}$ (at most $k-1$ of the first $n$ values lie below $V_{n+1}$, since all
values below $V_{n+1}$ lie below $Q$ and exactly $k-1$ of the first $n$ do); with distinct values
$\#\{j : R_j \ge k\} = n+1-k$ exactly, so $\Prob(R_{n+1} \le k-1) = 1 - (n+1-k)/(n+1) = k/(n+1) \le
1-\alpha+1/(n+1)$.
\end{proof}

\paragraph{Claim (Theorem~\ref{thm:xgb-cqr}).}
\emph{In split CQR with calibration set $I_2$ of size $n_2$, if the calibration pairs $\{(X_j,
Y_j)\}_{j \in I_2}$ and the test pair $(X, Y)$ are exchangeable, then $\Prob\{Y \in C(X)\} \ge 1 -
\alpha$; if the conformity scores are almost surely distinct, also $\Prob\{Y \in C(X)\} \le 1 -
\alpha + 1/(n_2+1)$ \cite{romano2019}.}

\begin{proof}
Condition on the proper training set $I_1$; the fitted quantile functions $\hat q_{\alpha/2}, \hat
q_{1-\alpha/2}$ are then fixed measurable functions. Define the score map $s(x,y) = \max\{\hat
q_{\alpha/2}(x) - y,\; y - \hat q_{1-\alpha/2}(x)\}$ and set $E_j = s(X_j, Y_j)$ for $j \in I_2$ and
$E_{n_2+1} = s(X, Y)$. A fixed measurable function applied coordinatewise to an exchangeable vector
preserves exchangeability, so $(E_j)_{j \in I_2}$ together with $E_{n_2+1}$ are exchangeable given
$I_1$. The algorithm's threshold $Q_{1-\alpha}$ is the $k$-th smallest calibration score with $k =
\lceil (n_2+1)(1-\alpha) \rceil$ (if $k > n_2$, set $Q_{1-\alpha} = +\infty$ and the coverage bound
is trivial). The coverage event is exactly a score comparison:
\[
Y \in C(X)
\;\Longleftrightarrow\;
\hat q_{\alpha/2}(X) - Y \le Q_{1-\alpha} \ \text{and}\ Y - \hat q_{1-\alpha/2}(X) \le Q_{1-\alpha}
\;\Longleftrightarrow\;
E_{n_2+1} \le Q_{1-\alpha}.
\]
Lemma~\ref{lem:app-rank}, applied conditionally on $I_1$, gives $\Prob\{Y \in C(X) \mid I_1\} \ge
1-\alpha$ and, under almost-sure distinctness, the matching upper bound. Integrating over $I_1$
preserves both inequalities. The finite-sample guarantee is \emph{marginal} and rests entirely on
exchangeability, which temporal drift breaks --- Warning~\ref{warn:xgb-exch} is the operational
consequence, not a footnote.
\end{proof}

\section{Proofs for Chapter~\ref{ch:causal}}
\label{app:proofs-ch5}

We work with discrete variables and a strictly positive joint law; the continuous case is identical
with densities. Let $G$ be a DAG over $V = (V_1, \dots, V_p)$ and let $P$ factorize as $P(v) =
\prod_j P(v_j \mid \mathrm{pa}_j)$. Following \cite{pearl1995}, the interventional law is
\emph{defined} by the truncated factorization
\begin{equation}
\label{eq:app-trunc}
P_a(v) \;=\; \prod_{j : V_j \ne A} P(v_j \mid \mathrm{pa}_j) \Big|_{A = a}
\qquad \text{on } \{v : v_A = a\},
\end{equation}
and $\E[Y \mid \Do(A=a)]$ is the mean of $Y$ under $P_a$.

\begin{lemma}[Interventions do not move non-descendants]
\label{lem:app-truncmarg}
Let $W$ be the set of non-descendants of $A$ (excluding $A$). Then $P_a(w) = P(w)$. In particular
$P_a(z) = P(z)$ for any $Z \subseteq W$.
\end{lemma}

\begin{proof}
No parent of a $W$-variable can be $A$ or a strict descendant $D$-variable of $A$: a child of $A$ is
a descendant, and a variable with a parent in $D$ is itself a descendant. Likewise every parent of
$A$ is in $W$ (a parent that is a descendant would close a directed cycle). Hence the order ($W$ in
topological order, then $A$, then the strict descendants $D$ in topological order) is a topological
order of $G$, and
\[
P(w, a, d) = \Big[\textstyle\prod_{j \in W} P(w_j \mid \mathrm{pa}_j)\Big]\, P(a \mid \mathrm{pa}_A)
\Big[\textstyle\prod_{j \in D} P(d_j \mid \mathrm{pa}_j)\Big],
\qquad
P_a(w, d) = \Big[\textstyle\prod_{j \in W}\Big]\Big[\textstyle\prod_{j \in D}\Big]\Big|_{A=a}.
\]
Summing $P_a(w,d)$ over the $D$-coordinates in reverse topological order telescopes (each
conditional sums to one), leaving $P_a(w) = \prod_{j \in W} P(w_j \mid \mathrm{pa}_j)$. The same
telescoping in the observational law (sum over $D$, then over $a$) gives $P(w) = \prod_{j \in W}
P(w_j \mid \mathrm{pa}_j)$. Summing out $W \setminus Z$ gives the claim for $Z$, which contains no
descendant of $A$ by backdoor condition (i).
\end{proof}

\begin{lemma}[Backdoor blocking survives edge deletion]
\label{lem:app-block}
If $Z$ satisfies the backdoor criterion (Definition~\ref{def:causal-backdoor}) for $(A, Y)$ in $G$,
then $Z$ d-separates $A$ from $Y$ in the graph $G_{\underline{A}}$ obtained from $G$ by deleting all
edges out of $A$.
\end{lemma}

\begin{proof}
Let $\pi$ be any $A$--$Y$ path in $G_{\underline{A}}$. Its first edge points into $A$ (all out-edges
are deleted), so $\pi$ is also a path of $G$ and is a backdoor path there; by criterion (ii) it is
blocked by $Z$ in $G$: some node $W$ on $\pi$ is either (a) a non-collider on $\pi$ with $W \in Z$,
or (b) a collider on $\pi$ with neither $W$ nor any $G$-descendant of $W$ in $Z$. Collider status
depends only on the orientations of the two $\pi$-edges at $W$, which are the same in both graphs.
In case (a), $\pi$ is blocked at $W$ in $G_{\underline{A}}$ as well. In case (b), the
$G_{\underline{A}}$-descendants of $W$ form a subset of its $G$-descendants, so again neither $W$
nor any $G_{\underline{A}}$-descendant of $W$ lies in $Z$, and $\pi$ is blocked. Hence every
$A$--$Y$ path of $G_{\underline{A}}$ is blocked by $Z$.
\end{proof}

\paragraph{Claim (Theorem~\ref{thm:causal-adjustment}).}
\emph{If $Z$ satisfies the backdoor criterion relative to $(A, Y)$, then $\E[Y \mid \Do(A=a)] =
\E_Z[\, \E[Y \mid A=a, Z] \,]$, and consequently $\tau_{\mathrm{ATE}} = \E_Z[\E[Y \mid A{=}1, Z] -
\E[Y \mid A{=}0, Z]]$.}

\begin{proof}
Decompose $P_a(y) = \sum_z P_a(y \mid z)\, P_a(z)$. Lemma~\ref{lem:app-truncmarg} gives $P_a(z) =
P(z)$. For the conditional we invoke the invariance half of Rule~2 of the do-calculus: \emph{if $Z$
d-separates $A$ from $Y$ in $G_{\underline{A}}$, then $P_a(y \mid z) = P(y \mid a, z)$}. This is
proved in \cite{pearl1995} from the truncated factorization \eqref{eq:app-trunc} via the soundness
of d-separation for DAG-factorized laws; we delegate exactly this step, and
Lemma~\ref{lem:app-block} verifies its graphical hypothesis from the backdoor definition. Combining,
$P_a(y) = \sum_z P(y \mid a, z) P(z)$, i.e.\ $\E[Y \mid \Do(A=a)] = \sum_z P(z)\, \E[Y \mid A=a,
Z=z]$; subtracting the $a = 0$ and $a = 1$ versions yields the ATE formula.

Two special cases are worth recording because they are fully self-contained. (1) If $Z =
\mathrm{pa}_A$, no graphical lemma is needed: on $\{A = a\}$, $P(v) = P_a(v)\, P(a \mid z)$ directly
from \eqref{eq:app-trunc}, so $P_a(y, z) = P(y, z, a)/P(a \mid z) = P(y \mid a, z) P(z)$, which is
the formula. (2) Under Assumption~\ref{ass:causal-ignorability} with potential-outcome semantics
($\E[Y \mid \Do(A=a)] := \E[Y(a)]$), the formula is three lines: $\E[Y(a)] = \E_Z\E[Y(a) \mid Z] =
\E_Z\E[Y(a) \mid A=a, Z] = \E_Z\E[Y \mid A=a, Z]$, using ignorability and then consistency
(Assumption~\ref{ass:causal-consistency}). Case (2) is the form under which Chapter~\ref{ch:causal}
actually operates.
\end{proof}

\paragraph{Claim (Proposition~\ref{prop:causal-dr}).}
\emph{Let $\psi$ be the AIPW score \eqref{eq:causal-aipw} built from working nuisances $(\mu_1,
\mu_0, e)$ with $\delta \le e(Z) \le 1 - \delta$ for some $\delta > 0$ and $\E|Y| + \E|\mu_a(Z)| <
\infty$. Under Assumptions~\ref{ass:causal-consistency}--\ref{ass:causal-ignorability}, $\E[\psi] =
\tau_{\mathrm{ATE}}$ if either $(\mu_1, \mu_0) = (\mu_1^*, \mu_0^*)$ is the true outcome pair or $e
= e^*$ is the true propensity.}

\begin{proof}
Write $\psi = \psi_1 - \psi_0$ with $\psi_1 = \mu_1(Z) + A\,(Y - \mu_1(Z))/e(Z)$ and $\psi_0 =
\mu_0(Z) + (1-A)(Y - \mu_0(Z))/(1 - e(Z))$; all terms are integrable because the denominators are
bounded away from zero. We show $\E[\psi_1] = \E[Y(1)]$ under either hypothesis; the $\psi_0$ half
is symmetric, and subtracting gives the claim.

\emph{True outcome model.} Recall $\mu_1^*(Z) = \E[Y \mid A=1, Z] = \E[Y(1) \mid A=1, Z] = \E[Y(1)
\mid Z]$, by consistency and then ignorability. Conditioning on $Z$ and using that $e(Z)$ is
$Z$-measurable,
\[
\E\Big[\frac{A\,(Y - \mu_1^*(Z))}{e(Z)} \,\Big|\, Z\Big]
= \frac{\Prob(A{=}1 \mid Z)}{e(Z)}\;\E[Y - \mu_1^*(Z) \mid A=1, Z]
= \frac{e^*(Z)}{e(Z)} \cdot 0 = 0 ,
\]
regardless of the working $e$. Hence $\E[\psi_1] = \E[\mu_1^*(Z)] = \E[\E[Y(1) \mid Z]] = \E[Y(1)]$.

\emph{True propensity.} With $e = e^*$ and arbitrary $\mu_1$,
\[
\E[\psi_1 \mid Z]
= \mu_1(Z) + \frac{e^*(Z)\,\E[Y \mid A=1, Z] - e^*(Z)\,\mu_1(Z)}{e^*(Z)}
= \E[Y \mid A=1, Z] = \E[Y(1) \mid Z],
\]
the working $\mu_1$ cancelling exactly; take expectations over $Z$.
\end{proof}

\begin{proposition}[Influence-function variance of AIPW]
\label{prop:app-aipwvar}
Let both nuisances be correct, $\E[Y^2] < \infty$, and let $O_1, \dots, O_n$ be i.i.d. Write
$\tau(Z) = \mu_1^*(Z) - \mu_0^*(Z)$ and $\sigma_a^2(Z) = \Var(Y \mid A=a, Z)$. Then
$\hat\tau_{\mathrm{AIPW}} = n^{-1}\sum_i \psi(O_i)$ satisfies $\sqrt{n}\,(\hat\tau_{\mathrm{AIPW}} -
\tau_{\mathrm{ATE}}) \to \mathcal{N}(0, V)$ with
\[
V \;=\; \Var(\psi) \;=\; \Var\big(\tau(Z)\big)
\;+\; \E\Big[\frac{\sigma_1^2(Z)}{e^*(Z)}\Big]
\;+\; \E\Big[\frac{\sigma_0^2(Z)}{1 - e^*(Z)}\Big].
\]
\end{proposition}

\begin{proof}
Asymptotic normality is the Lindeberg--Levy central limit theorem for the i.i.d.\ mean of the
square-integrable $\psi$, whose mean is $\tau_{\mathrm{ATE}}$ by Proposition~\ref{prop:causal-dr}.
For the variance, write $\psi = \tau(Z) + U_1 - U_0$ with $U_1 = A(Y - \mu_1^*(Z))/e^*(Z)$ and $U_0
= (1-A)(Y - \mu_0^*(Z))/(1 - e^*(Z))$. Both $U_a$ have conditional mean zero given $Z$ (shown in the
previous proof), so they are uncorrelated with the $Z$-measurable $\tau(Z)$, and $U_1 U_0 = 0$
pointwise because $A(1-A) = 0$. By conditioning on $(A, Z)$ and using $\mu_1^*(Z) = \E[Y \mid A=1,
Z]$,
\[
\E[U_1^2 \mid Z] = \frac{\E[A (Y - \mu_1^*(Z))^2 \mid Z]}{e^*(Z)^2}
= \frac{e^*(Z)\, \sigma_1^2(Z)}{e^*(Z)^2} = \frac{\sigma_1^2(Z)}{e^*(Z)},
\]
and symmetrically for $U_0$. Adding the three orthogonal contributions gives $V$. This $V$ is the
semiparametric efficiency bound for the ATE, and the cross-fitted estimator of
Section~\ref{sec:causal-estimators} attains it under the product-rate condition of
Remark~\ref{rem:causal-rates}; see \cite{chernozhukov2018}. The firm-clustered variance
\eqref{eq:causal-clustervar} is the sample analogue of $\Var(\sum_y \psi_{i,y})$ with firms as
independent blocks, and the cluster bootstrap of Section~\ref{sec:causal-inference} estimates the
same quantity by resampling \cite{efron1994}.
\end{proof}

\section{Proofs for Chapter~\ref{ch:us}}
\label{app:proofs-ch6}

Chapter~\ref{ch:us} uses three display-level identities --- the selection-bias covariance formula
\eqref{eq:us-selbias}, the worst-case bounds \eqref{eq:us-bounds}, and the unbiasedness of the
inverse-probability-of-observation weights below \eqref{eq:us-pi}. We record them with proofs.

\begin{proposition}[Selection identities]
\label{prop:app-selection}
Let $g$ be integrable, $D \in \{0,1\}$ with $p = \Prob(D{=}1) \in (0,1)$. Then: (i) $\E[g \mid
D{=}1] - \E[g] = \Cov(g, \ind{D{=}1})/p$; (ii) if $g \in [g_{\mathrm{lo}}, g_{\mathrm{hi}}]$ almost
surely, then $\E[g]$ lies in the interval \eqref{eq:us-bounds}; (iii) under
Assumption~\ref{ass:us-soh}, with $\pi = \Prob(D{=}1 \mid \Hist_{T_{i,k}}) > 0$ almost surely, $\E[D
g/\pi] = \E[g]$.
\end{proposition}

\begin{proof}
(i) $\Cov(g, \ind{D{=}1}) = \E[g \ind{D{=}1}] - p\,\E[g] = p\,\E[g \mid D{=}1] - p\,\E[g]$; divide
by $p$. (ii) By the law of total expectation, $\E[g] = p\,\E[g \mid D{=}1] + (1-p)\,\E[g \mid
D{=}0]$, and $\E[g \mid D{=}0] \in [g_{\mathrm{lo}}, g_{\mathrm{hi}}]$ by the almost-sure bounds.
(iii) By the tower property and the conditional independence $g \indep D \mid \Hist_{T_{i,k}}$,
\[
\E\Big[\frac{D g}{\pi}\Big]
= \E\Big[\frac{1}{\pi}\, \E[D g \mid \Hist_{T_{i,k}}]\Big]
= \E\Big[\frac{1}{\pi}\, \E[D \mid \Hist_{T_{i,k}}]\; \E[g \mid \Hist_{T_{i,k}}]\Big]
= \E\big[\E[g \mid \Hist_{T_{i,k}}]\big] = \E[g],
\]
all steps valid because $\pi$ is $\Hist_{T_{i,k}}$-measurable and positive.
\end{proof}

\begin{remark}
\label{rem:app-pi}
The closed form \eqref{eq:us-pi} for $\pi$ is the survival function of the first event after
$T_{i,k}$: with an internal (self-exciting) history, the intensity along the no-new-event
continuation is an $\Hist_{T_{i,k}}$-measurable function of elapsed time, and the hazard
representation of the inter-event distribution \cite[Ch.~7.2]{daley2003} gives $\Prob(T_{i,k+1} >
\tau_{\mathrm{end}} \mid \Hist_{T_{i,k}}) = \exp(-\int \lambda_i)$. If external covariates evolve
stochastically over the gap, \eqref{eq:us-pi} holds with the exponential replaced by its conditional
expectation.
\end{remark}

\section{Proofs for Chapter~\ref{ch:poisson}}
\label{app:proofs-ch7}

\paragraph{Claim (Theorem~\ref{thm:poi-watanabe}).}
\emph{Let $N$ be a simple point process adapted to $\Hist_t$ with continuous compensator $\Lambda$.
Then $N$ is an inhomogeneous Poisson process --- independent increments, $N(b) - N(a) \sim
\mathrm{Poisson}(\Lambda(b) - \Lambda(a))$ --- if and only if $\Lambda$ is deterministic.}

\begin{proof}
\emph{Deterministic implies Poisson (the direction the manual uses).} Assume $\Lambda$ is a
continuous deterministic nondecreasing function, finite on compacts, with $N - \Lambda$ a martingale
(Lemma~\ref{lem:app-mart}(i); here \eqref{eq:app-integrable} holds automatically since $\Lambda$ is
finite). Fix $u \in \mathbb{R}$ and $a \ge 0$, and define for $t \ge a$
\begin{equation}
\label{eq:app-expmart}
Z_t \;=\; \exp\Big\{ i u \big(N(t) - N(a)\big) - \big(e^{iu} - 1\big)\big(\Lambda(t) -
\Lambda(a)\big) \Big\} \;=\; e^{iu(N(t)-N(a))}\, g(t),
\end{equation}
with $g$ continuous, deterministic, of bounded variation, and bounded away from $0$ and $\infty$ on
compacts. Both factors are of bounded variation, so Stieltjes integration by parts (with $g$
continuous) gives $\dd Z_t = e^{iu(N(t^-)-N(a))}\, \dd g(t) + g(t)\, \dd\big(e^{iu(N-N(a))}\big)_t$.
The first differential is $-\,(e^{iu}-1)\, Z_{t^-}\, \dd\Lambda(t)$ by the chain rule for continuous
BV functions (and $Z_{t^-}$ may replace $Z_t$ in a $\dd\Lambda$-integral because $\Lambda$, being
continuous, charges no countable set, in particular not the jump times). The second is a pure-jump
differential: at an event, $e^{iuN}$ is multiplied by $e^{iu}$, so $\dd\big(e^{iu(N - N(a))}\big)_t
= (e^{iu}-1)\, e^{iu(N(t^-)-N(a))}\, \dd N(t)$. Adding,
\[
Z_t \;=\; Z_a + \int_{(a,t]} Z_{s^-}\,(e^{iu}-1)\; \dd\big(N - \Lambda\big)(s).
\]
The integrand $H_s = Z_{s^-}(e^{iu}-1)$ is predictable (left limit of an adapted cadlag process) and
bounded on $[a,b]$ by $|e^{iu}-1|\, e^{|e^{iu}-1|\Lambda(b)}$, a deterministic constant. By
Lemma~\ref{lem:app-mart}(ii) the integral is a martingale, so $\E[Z_t \mid \Hist_a] = Z_a = 1$, i.e.
\[
\E\big[ e^{iu(N(t) - N(a))} \,\big|\, \Hist_a \big]
\;=\; \exp\big\{ (e^{iu} - 1)(\Lambda(t) - \Lambda(a)) \big\},
\]
the characteristic function of $\mathrm{Poisson}(\Lambda(t)-\Lambda(a))$, and it is
\emph{deterministic}: by uniqueness of characteristic functions the increment is Poisson
distributed, and since the conditional characteristic function is constant, the increment is
independent of $\Hist_a$. For disjoint intervals $(a_1,b_1], \dots, (a_m,b_m]$ (ordered), iterating
the tower property from the last interval inward factorizes the joint characteristic function of the
increments into the product of Poisson characteristic functions, proving joint independence. Hence
$N$ is inhomogeneous Poisson with mean function $\Lambda$.

\emph{Poisson implies deterministic.} If $N$ has independent Poisson increments with continuous mean
function $m$, then $N(t) - m(t)$ is a martingale ($\E[N(t) - N(s) \mid \Hist_s] = m(t)-m(s)$ by
independence), and $m$ is a continuous deterministic --- hence predictable --- increasing process.
The compensator is the unique predictable increasing process making $N$ minus it a martingale
(Doob--Meyer uniqueness, which we take from \cite{bremaud1996}), so $\Lambda = m$ almost surely:
deterministic. For the general marked-space version of both directions see \cite[Ch.~7]{daley2003}.
\end{proof}

\paragraph{Claim (Theorem~\ref{thm:poi-rescaling}).}
\emph{Let $N$ have event times $t_1 < t_2 < \cdots$ and conditional intensity $\lambda(t) > 0$ with
$\Lambda(t) \to \infty$ almost surely. Then $\tau_k = \Lambda(t_k)$ are the event times of a
unit-rate Poisson process; equivalently $\delta_k = \tau_k - \tau_{k-1}$ are i.i.d.\
$\mathrm{Exp}(1)$.}

\begin{proof}
\emph{Step 1: the time change.} $\Lambda$ is continuous (an integral) and strictly increasing: for
$a<b$, $\Lambda(b)-\Lambda(a) = \int_a^b \lambda > 0$ because $\lambda$ is strictly positive. Since
$\Lambda(\infty) = \infty$ a.s., for each $s \ge 0$ there is a unique finite $\tau(s)$ with
$\Lambda(\tau(s)) = s$, and $s \mapsto \tau(s)$ is continuous and increasing. Each $\tau(s)$ is a
stopping time: $\{\tau(s) \le t\} = \{\Lambda(t) \ge s\} \in \Hist_t$. Define $\mathcal{G}_s =
\Hist_{\tau(s)}$ (a filtration, since $\tau$ is monotone) and $\tilde N(s) = N(\tau(s))$. Because
$N$ is right-continuous and adapted, it is progressively measurable, so $\tilde N(s)$ is
$\mathcal{G}_s$-measurable (the standard measurability of a progressive process at a stopping time).
$\tilde N$ is a simple counting path whose jump times are exactly $\{\Lambda(t_k)\} = \{\tau_k\}$,
since the strictly increasing $\Lambda$ maps distinct event times to distinct values.

\emph{Step 2: integrability.} Fix $s$ and let $T = \tau(s)$. For each $n$, $T \wedge n$ is a bounded
stopping time, and Doob's optional sampling applied to the martingale $M = N - \Lambda$ on $[0,n]$
gives $\E[N(T \wedge n)] = \E[\Lambda(T \wedge n)] \le s$. Monotone convergence ($T < \infty$ a.s.,
$N$ and $\Lambda$ nondecreasing) yields $\E[\tilde N(s)] = \E[\Lambda(T)] = s < \infty$.

\emph{Step 3: $\tilde N(s) - s$ is a $\mathcal{G}$-martingale.} Fix $s' < s$, put $T' = \tau(s')$,
$T = \tau(s)$, and take $A \in \mathcal{G}_{s'} = \Hist_{T'}$. For each $n$, $T' \wedge n \le T
\wedge n$ are bounded stopping times, so optional sampling gives $\E[M(T \wedge n) \mid \Hist_{T'
\wedge n}] = M(T' \wedge n)$. The event $A_n = A \cap \{T' \le n\}$ belongs to $\Hist_{T' \wedge n}$
(for $t < n$, $A_n \cap \{T' \wedge n \le t\} = A \cap \{T' \le t\} \in \Hist_t$; for $t \ge n$ it
lies in $\Hist_n \subseteq \Hist_t$), whence $\E[\ind{A_n} M(T \wedge n)] = \E[\ind{A_n} M(T' \wedge
n)]$. As $n \to \infty$: $\ind{A_n} \uparrow \ind{A}$ since $T' < \infty$ a.s.; $M(T \wedge n) \to
M(T)$ and $M(T' \wedge n) \to M(T')$ a.s.; and $|M(T \wedge n)| \le N(T) + s$ with $N(T) \in L^1$ by
Step 2 (similarly for $T'$), so dominated convergence applies on both sides: $\E[\ind{A} M(T)] =
\E[\ind{A} M(T')]$ for all $A \in \mathcal{G}_{s'}$. Since $M(T') = \tilde N(s') - s'$ is
$\mathcal{G}_{s'}$-measurable, this says $\E[\tilde N(s) - s \mid \mathcal{G}_{s'}] = \tilde N(s') -
s'$.

\emph{Step 4: conclude.} $\tilde N$ is a simple counting process adapted to $(\mathcal{G}_s)$ whose
compensator is the continuous deterministic function $s$. By the direction of
Theorem~\ref{thm:poi-watanabe} proved above (whose proof used only the martingale property of $N -
\Lambda$ with $\Lambda$ continuous deterministic), $\tilde N$ is a unit-rate Poisson process. Its
jump times are the $\tau_k$, and the equivalence of a unit-rate Poisson process with i.i.d.\
standard-exponential interarrivals is the classical description of the Poisson process
\cite[Ch.~2]{daley2003}; the residual-analysis use is due to \cite{ogata1988}.
\end{proof}

\paragraph{Claim (Corollary~\ref{cor:poi-testable}).}
\emph{Under correct specification, every fitted model class in Part~III yields rescaled
inter-arrivals $\hat\delta_k$ that are i.i.d.\ $\mathrm{Exp}(1)$, and $u_k = 1 - e^{-\hat\delta_k}$
that are i.i.d.\ $\mathrm{Uniform}(0,1)$.}

\begin{proof}
Every model class specifies $\lambda(t; \theta)$ explicitly with $\lambda > 0$; under correct
specification the data are generated by $\lambda(\cdot\,; \theta_0)$, and
Theorem~\ref{thm:poi-rescaling} applied at $\theta_0$ makes the rescaled inter-arrivals i.i.d.\
$\mathrm{Exp}(1)$. The probability integral transform finishes: $\Prob(1 - e^{-\delta} \le u) =
\Prob(\delta \le -\log(1-u)) = u$ for $u \in (0,1)$, and measurable transformations of independent
variables are independent. (In practice $\hat\theta$ replaces $\theta_0$; the tests are then
conservative screens, exact only asymptotically.)
\end{proof}

\paragraph{Claim (Proposition~\ref{prop:poi-concave}).}
\emph{With $x_w = (1, X_w^\top)^\top$, the weekly log-likelihood \eqref{eq:poi-loglik} is concave in
$\gamma = (a, \beta)$, strictly concave iff $\{x_w\}$ spans $\mathbb{R}^{p+1}$; the MLE, when it
exists, is then unique.}

\begin{proof}
Write $\eta_w = x_w^\top \gamma$, so $\ell(\gamma) = \sum_w [\, n_w \eta_w - \Delta_w e^{\eta_w}
\,]$. Then $\nabla \ell = \sum_w (n_w - \Delta_w e^{\eta_w})\, x_w$ and $\nabla^2 \ell = -\sum_w
\Delta_w e^{\eta_w} x_w x_w^\top$. For any $v$, $v^\top \nabla^2 \ell\, v = -\sum_w \Delta_w
e^{\eta_w} (x_w^\top v)^2 \le 0$, so $\ell$ is concave; the form is zero iff $x_w^\top v = 0$ for
all $w$. If the $x_w$ span, no such $v \ne 0$ exists and $\nabla^2 \ell \prec 0$ everywhere: $\ell$
is strictly concave, and two distinct maximizers would contradict strict concavity at their
midpoint, so the maximizer is unique. Conversely, if some $v \ne 0$ is orthogonal to every $x_w$,
then $\ell(\gamma + t v) = \ell(\gamma)$ for all $t$: $\ell$ is constant along a line, hence not
strictly concave and its maximizer (if any) not unique.
\end{proof}

\paragraph{Claim (Proposition~\ref{prop:poi-ident}).}
\emph{The Fisher information of \eqref{eq:poi-loglik} is $I(a, \beta) = \sum_w \Delta_w \lambda_w
x_w x_w^\top$ with $\lambda_w = e^{a + \beta^\top X_w}$, and after profiling out $a$ the information
for $\beta$ is the intensity-weighted temporal covariance \eqref{eq:poi-fisher}. A covariate
constant over the window is not identified separately from $a$; a near-linear trend leaves as
effective information the weighted variance of its detrended residual.}

\begin{proof}
The Hessian computed in the previous proof does not involve the counts $n_w$, so observed and
expected information coincide: $I(\gamma) = -\nabla^2\ell = \sum_w \omega_w x_w x_w^\top$ with
$\omega_w = \Delta_w \lambda_w$. Partition by the intercept: $I_{aa} = \sum_w \omega_w$, $I_{a\beta}
= \sum_w \omega_w X_w^\top$, $I_{\beta\beta} = \sum_w \omega_w X_w X_w^\top$. The information for
$\beta$ with $a$ unknown is the Schur complement $I_\beta = I_{\beta\beta} - I_{\beta a} I_{aa}^{-1}
I_{a\beta}$ (the inverse of $[I^{-1}]_{\beta\beta}$, by the block-inverse formula). With $\bar
X_\omega = \sum_w \omega_w X_w / \sum_w \omega_w$,
\[
\sum_w \omega_w (X_w - \bar X_\omega)(X_w - \bar X_\omega)^\top
= \sum_w \omega_w X_w X_w^\top - \Big(\sum_w \omega_w\Big) \bar X_\omega \bar X_\omega^\top
= I_{\beta\beta} - I_{\beta a} I_{aa}^{-1} I_{a\beta},
\]
which is \eqref{eq:poi-fisher}. (i) If coordinate $j$ of $X_w$ is constant, then $X_w^{(j)} - \bar
X_\omega^{(j)} = 0$ for all $w$, so $I_\beta$ has a zero row and column: singular. The
non-identification is exact, not asymptotic: replacing $(a, \beta_j) \to (a - tc, \beta_j + t)$
leaves every linear predictor $\eta_w$, hence the entire likelihood, unchanged. (ii) More generally,
let $U$ collect nuisance columns (intercept and an explicit trend $t_w$) and endow $\mathbb{R}^W$
with the inner product $\langle f, g \rangle_\omega = \sum_w \omega_w f_w g_w$. The Schur complement
of the $U$-block in the full Gram matrix is the Gram matrix of the residuals of the $X$-columns
after $\langle\cdot,\cdot\rangle_\omega$-orthogonal projection onto $\mathrm{span}(U)$ --- because
$I_{\beta U} I_{UU}^{-1} I_{U \beta}$ is precisely the Gram matrix of those projections. Hence the
profile information for $\beta$ is $\sum_w \omega_w r_w r_w^\top$ with $r_w$ the $\omega$-weighted
regression residual of $X_w$ on $(1, t_w)$: a covariate that is nearly a linear trend has small
residuals and correspondingly little information, as claimed.
\end{proof}

\begin{proposition}[Instability of log-link raw-count feedback (Mechanism B)]
\label{prop:app-loglink}
Let $a \in \mathbb{R}$, $b > 0$, and let $Y_t \mid \Hist_t \sim \mathrm{Poisson}(\lambda_t)$ with
$\lambda_{t+1} = \exp(a + b Y_t)$ and any initial $\lambda_1 > 0$. Define $G(\lambda) = e^{a}
\exp\{\lambda (e^{b} - 1)\}$. Then: (i) $\E[\lambda_{t+1} \mid \Hist_t] = G(\lambda_t)$ exactly,
with $G$ increasing, strictly convex, and $G'(\lambda^*) = \lambda^* (e^b - 1)$ at any fixed point
$\lambda^*$; (ii) $G$ has a fixed point iff $e^{a}(e^{b}-1) \le e^{-1}$, in which case there are at
most two, $\lambda^*_- \le \lambda^*_+$, with $\lambda^*_-(e^b - 1) \le 1$ (the local stability
condition of the mean recursion) and $\lambda^*_+(e^b - 1) \ge 1$ (repelling); (iii) in all cases
$\E[\lambda_t] \to \infty$ as $t \to \infty$: no fixed point is globally attracting for the mean
dynamics; (iv) by contrast, under linear feedback $\lambda_{t+1} = c + \rho Y_t$ with $c > 0$ and $0
\le \rho < 1$, $\E[\lambda_t] \to c/(1-\rho)$ geometrically from every initial condition.
\end{proposition}

\begin{proof}
(i) The Poisson moment generating function: for $Y \sim \mathrm{Poisson}(\lambda)$ and any real $s$,
$\E[e^{sY}] = \sum_{k \ge 0} e^{sk} e^{-\lambda}\lambda^k / k! = e^{\lambda(e^s - 1)}$ (the series
is the exponential series at $\lambda e^s$). Hence $\E[\lambda_{t+1} \mid \Hist_t] = e^a\,
\E[e^{bY_t} \mid \Hist_t] = G(\lambda_t)$. Direct differentiation gives $G' = (e^b - 1) G > 0$ and
$G'' = (e^b-1)^2 G > 0$; at a fixed point $G(\lambda^*) = \lambda^*$, so $G'(\lambda^*) =
\lambda^*(e^b - 1)$.

(ii) Write $c_0 = e^a$, $d = e^b - 1 > 0$. Fixed points solve $c_0 e^{d\lambda} = \lambda$. The
strictly convex function $g(\lambda) = G(\lambda) - \lambda$ has $g(0) = c_0 > 0$ and $g(\lambda)
\to \infty$; it attains its minimum where $G' = 1$, i.e.\ at $\tilde\lambda$ with $c_0 d\,
e^{d\tilde\lambda} = 1$, giving $\tilde\lambda = d^{-1}\log(1/(c_0 d))$ when positive and
$g(\tilde\lambda) = 1/d - \tilde\lambda$. Thus $g$ has a zero iff $\tilde\lambda \ge 1/d$, i.e.\ iff
$c_0 d \le e^{-1}$, which is the stated condition. A strictly convex function has at most two zeros;
at the smaller, $g$ crosses from positive to non-positive, so $g'(\lambda^*_-) \le 0$, i.e.\
$G'(\lambda^*_-) \le 1$; at the larger the crossing is upward, $G'(\lambda^*_+) \ge 1$.

(iii) Note $\lambda_t \ge e^a$ for every $t \ge 2$. Choose an integer $k$ with $\lambda' := e^{a +
bk}$ strictly above every fixed point (possible: fixed points are finite, or absent). Then
$G(\lambda) > \lambda$ for all $\lambda \ge \lambda'$, and the deterministic iterates $G^{\circ
n}(\lambda')$ increase; were they bounded, their limit would be a fixed point $\ge \lambda'$, a
contradiction --- so $G^{\circ n}(\lambda') \uparrow \infty$. Poisson tails are stochastically
increasing in the mean (superpose an independent $\mathrm{Poisson}(\mu' - \mu)$), so $p_0 :=
\Prob(\mathrm{Poisson}(e^a) \ge k) > 0$ lower-bounds $\Prob(Y_2 \ge k \mid \Hist_2)$. By conditional
Jensen with the convex increasing $G$ and induction on $n$, $\E[\lambda_{3+n} \mid \Hist_3] \ge
G^{\circ n}(\lambda_3)$ (allowing $+\infty$ throughout). On the event $B = \{Y_2 \ge k\}$ we have
$\lambda_3 = e^{a + bY_2} \ge \lambda'$, hence
\[
\E[\lambda_{3+n}] \;\ge\; \E\big[\ind{B}\, \E[\lambda_{3+n} \mid \Hist_3]\big]
\;\ge\; \Prob(B)\, G^{\circ n}(\lambda') \;\ge\; p_0\, G^{\circ n}(\lambda')
\;\longrightarrow\; \infty .
\]
(iv) Linearity passes through expectations exactly: $\E[\lambda_{t+1} \mid \Hist_t] = c + \rho
\lambda_t$, so $\E[\lambda_t] - c/(1-\rho) = \rho^{\,t-1}(\E[\lambda_1] - c/(1-\rho)) \to 0$. The
contrast with (iii) is precisely the Jensen gap: $G$ is convex, so rare large counts dominate the
mean, while the linear map has no gap. Note also $e^b - 1 > b$, so the true mean-map derivative
$\lambda^*(e^b - 1)$ strictly exceeds the naive effective radius $\rho_{\mathrm{eff}} = b\lambda^*$
of \eqref{eq:poi-rhoeff}: a configuration comfortably subcritical by $\rho_{\mathrm{eff}}$ can be
metastable, which is the measured episode of Warning~\ref{warn:poi-metastable} and
Table~\ref{tab:sel-stability}. Genuine ergodicity theory for log-linear count autoregressions (with
$\log(1+Y)$ feedback) is in \cite{fokianos2009}; for nonlinear Hawkes stability see
\cite{bremaud1996}.
\end{proof}

\section{Proofs for Chapter~\ref{ch:hawkes}}
\label{app:proofs-ch8}

\begin{lemma}[Branching moments and local finiteness]
\label{lem:app-branching}
Consider a cluster rooted at a point at time $v$: generation $0$ is the root; each point of
generation $n$ independently produces $\mathrm{Poisson}(\kappa)$ generation-$(n{+}1)$ points at lags
drawn i.i.d.\ from the density $\phi/\kappa$ on $(0,\infty)$. Let $Z_n$ be the size of generation
$n$ and $S_n$ a sum of $n$ i.i.d.\ lags. Then (i) $\E Z_n = \kappa^n$, and for $\kappa < 1$ the mean
total cluster size is $\sum_n \kappa^n = 1/(1-\kappa)$; (ii) the expected number of generation-$n$
points within $[v, v+T]$ is $\kappa^n \Prob(S_n \le T)$; (iii) for every $T$ and every $\kappa$,
$\sum_n \kappa^n \Prob(S_n \le T) < \infty$, so clusters are almost surely finite on compact
windows.
\end{lemma}

\begin{proof}
(i) Wald's identity for a Poisson number of nonnegative i.i.d.\ terms: if $K$ is independent of the
i.i.d.\ sequence $(X_i)$, conditioning on $K$ and applying Tonelli gives $\E \sum_{i \le K} X_i = \E
K \cdot \E X_1$. Applying it generation by generation, $\E Z_{n+1} = \E[\E[Z_{n+1} \mid \text{gen }
n]] = \kappa\, \E Z_n$, so $\E Z_n = \kappa^n$, and the total follows by Tonelli over $n$. (ii) Let
$m_n$ be the mean measure of generation-$n$ displacements from the root. Then $m_0 = \delta_0$ and,
by the same conditioning plus Tonelli, $m_{n+1}(B) = \int\!\!\int \ind{u \in B}\, \phi(u - p)\, \dd
u\, m_n(\dd p)$: the density of $m_n$ is $\kappa^n$ times the $n$-fold convolution of $\phi/\kappa$,
i.e.\ the law of $S_n$ scaled by $\kappa^n$; evaluate on $[0, T]$. (iii) Lags are a.s.\ strictly
positive, so by dominated convergence $\E e^{-\eta S_1} \downarrow 0$ as $\eta \uparrow \infty$;
choose $\eta$ with $r := \kappa\, \E e^{-\eta S_1} < 1$. Markov's inequality gives $\Prob(S_n \le T)
\le e^{\eta T} (\E e^{-\eta S_1})^n$, hence $\sum_n \kappa^n \Prob(S_n \le T) \le e^{\eta T}/(1 - r)
< \infty$. Finite expectation implies almost-sure finiteness.
\end{proof}

\paragraph{Claim (Theorem~\ref{thm:hw-cluster}).}
\emph{For $\phi \ge 0$ with $\kappa = \int \phi < 1$, the Hawkes process \eqref{eq:hw-intensity} is
equal in law to the Poisson cluster process (immigrants at rate $s(\cdot)$, i.i.d.\
Poisson($\kappa$) offspring at lag density $\phi/\kappa$, superposed over generations). Under
constant $s$ it admits a stationary version iff $\kappa < 1$, with mean intensity $s/(1-\kappa)$ and
mean cluster size $1/(1-\kappa)$; the multivariate statement holds with $\rho(K) < 1$, mean vector
$(I-K)^{-1}s$.}

\begin{proof}
\emph{Construction.} On a suitable product space let immigrants form a Poisson process with locally
integrable intensity $s(\cdot)$ on $[0,\infty)$, and attach to each immigrant an independent cluster
as in Lemma~\ref{lem:app-branching}. Let $N_c$ be the superposition of all points. By
Lemma~\ref{lem:app-branching}(ii)--(iii) and Campbell's formula over immigrants, $\E N_c[0,T] \le
\big(\int_0^T s\big) \sum_n \kappa^n \Prob(S_n \le T) < \infty$, so $N_c$ is almost surely locally
finite; it is simple because all lags and immigrant positions have nonatomic laws, so any two
prescribed points coincide with probability zero and there are countably many pairs.

\emph{Intensity with respect to the genealogical filtration.} Let $\mathcal{F}_t$ be generated by
all points born up to $t$ together with their genealogy. Decompose $N_c$ into the immigrant stream
and the direct-offspring streams $N^p$ of each born point $p$. By construction, given $p$, $N^p$ is
a Poisson process on $(p, \infty)$ with intensity $\phi(\cdot - p)$, independent of all other
building blocks; by independence of Poisson increments, $N^p(t) - \int_p^{t \vee p} \phi(u - p)\,
\dd u$ is an $\mathcal{F}$-martingale, as is the compensated immigrant stream. Partial sums over the
countably many streams are martingales, and they converge in $L^1$ for each $t$ (monotonically in
the counting and compensating parts separately, with finite limit expectation $\E N_c(t) < \infty$),
so the limit $N_c(t) - \int_0^t \lambda_c(u)\, \dd u$, with $\lambda_c(u) = s(u) + \sum_{p < u}
\phi(u - p)$, is an $\mathcal{F}$-martingale ($L^1$-limits pass through conditional expectations).

\emph{Reduction to the internal history.} $\lambda_c$ is exactly the Hawkes formula
\eqref{eq:hw-intensity} evaluated on the points of $N_c$: it depends on point positions only, not on
the genealogy, so it is predictable for the internal filtration $\Hist_t = \sigma(N_c(s), s \le t)$,
and $\Hist_t \subseteq \mathcal{F}_t$. The tower property then transfers the martingale property:
for $s < t$, $\E[\,N_c(t) - \smallint_0^t \lambda_c \mid \Hist_s] = \E[\,\E[\cdots \mid
\mathcal{F}_s] \mid \Hist_s] = \E[\,N_c(s) - \smallint_0^s \lambda_c \mid \Hist_s] = N_c(s) -
\smallint_0^s \lambda_c$, the last equality because the variable is $\Hist_s$-measurable. Hence
$N_c$ has internal-history intensity \eqref{eq:hw-intensity}.

\emph{Law identification.} A simple point process whose internal-history intensity coincides almost
surely with a fixed predictable functional of the past configuration has its law on every $[0,T]$
uniquely determined --- the intensity determines the successive conditional hazard functions of
inter-event times, hence all finite-dimensional joint densities. We take this uniqueness from
\cite[Ch.~7.2--7.3]{daley2003}. Therefore $N_c$ equals in law the Hawkes process of
Definition~\ref{def:hw-hawkes}; this is the cluster representation, originally implicit in
\cite{hawkes1971}.

\emph{Stationarity.} Let $s$ be constant. For $\kappa < 1$, run the same construction with
immigrants homogeneous Poisson on all of $\mathbb{R}$. The construction mechanism is invariant under
time shifts, so the law of $N_c$ is stationary, and by Wald plus Campbell the mean rate is $s \cdot
\E[\text{cluster size}] = s/(1-\kappa) < \infty$ (Lemma~\ref{lem:app-branching}(i)). Conversely,
suppose a stationary version with finite mean intensity $\bar\xi = \E N(0,1]$ exists for some
$\kappa$. Then $\E \lambda(t)$ is finite and constant in $t$ (stationarity; and $\E \int_0^1 \lambda
= \E N(0,1]$ by Lemma~\ref{lem:app-fubini} with $f \equiv 1$), so $\E\lambda(t) = \bar\xi$ for a.e.\
$t$. Taking expectations in \eqref{eq:hw-intensity} and using Campbell's formula for the stationary
process ($\E \sum_{t_k < t} \phi(t - t_k) = \bar\xi \int_0^\infty \phi$, \cite{daley2003}) gives
$\bar\xi = s + \kappa \bar\xi$, i.e.\ $(1-\kappa)\bar\xi = s$. For $s > 0$ this has no finite
nonnegative solution when $\kappa \ge 1$; when $\kappa < 1$ it returns $\bar\xi = s/(1-\kappa)$,
consistent with the construction. The multivariate case repeats every step with matrix kernels:
generation means are $K^n$, $\sum_n K^n = (I - K)^{-1}$ converges iff $\rho(K) < 1$ (Gelfand's
formula), the mean vector is $(I-K)^{-1}s$, and the first-moment identity reads $(I - K)\bar\xi =
s$. For stability theory beyond the linear case see \cite{bremaud1996}.
\end{proof}

\paragraph{Claim (Proposition~\ref{prop:hw-nonneg}).}
\emph{Let $\lambda(t) = s(t) + \sum_{t_k < t} \phi(t - t_k)$ with $s$ bounded. If $\phi < 0$ on a
set of positive Lebesgue measure, there exist finite histories on which the formula is strictly
negative; hence \eqref{eq:hw-intensity} defines a valid intensity for all histories only if $\phi
\ge 0$ almost everywhere.}

\begin{proof}
Suppose $\phi$ is not nonnegative a.e. Then there are $\epsilon > 0$ and a Borel set $A \subseteq
(0, \infty)$ of positive Lebesgue measure with $\phi < -\epsilon$ on $A$. Let $s_{\max} = \sup_t
s(t) < \infty$ and fix an integer $n > s_{\max}/\epsilon$. A set of positive measure is infinite, so
we may choose $n$ distinct lags $a_1, \dots, a_n \in A$; fix any $t$ larger than all of them and
place events at $x_i = t - a_i$, a simple finite configuration. Evaluating the formula at $t$:
$\lambda(t) \le s_{\max} + \sum_{i=1}^n \phi(a_i) < s_{\max} - n\epsilon < 0$. A conditional
intensity is by definition nonnegative, so no point process can have \eqref{eq:hw-intensity} as its
intensity functional on \emph{all} finite histories unless $\phi \ge 0$ a.e.
\end{proof}

\begin{remark}[Why the all-histories qualifier is essential]
\label{rem:app-signed}
Requiring only \emph{almost-sure} nonnegativity along realized histories does not force $\phi \ge
0$: with constant baseline $s$ and $\phi = -s\,\ind{0 < u \le 1}$, the induced dynamics are a
renewal process with a unit dead time --- after each event the formula sits at exactly zero for one
time unit, so no two events are ever closer than $1$ and the offending configurations of the proof
are never visited. Such exactly tuned signed kernels are fragile (any perturbation of $s$ or $\phi$
breaks the cancellation), evade none of the substantive losses --- the cluster representation of
Theorem~\ref{thm:hw-cluster} requires $\phi \ge 0$ to define offspring densities, and truncating at
zero, $\lambda = (s + \sum\phi)_+$, destroys the linearity on which the Volterra first-moment
equation of Chapter~\ref{ch:censored} rests --- and none of the estimation theory of this manual
applies to them. The factorization route (b) of Section~\ref{sec:hw-inhib} avoids the problem
wholesale.
\end{remark}

\paragraph{Claim (Theorem~\ref{thm:hw-factor}).}
\emph{Let $N_s$ be the ground process of sector-$s$ deals with intensity $\lambda^g_s(t)$ and let
each event carry a firm mark. Then $\lambda_i(t) = \lambda^g_{s(i)}(t)\, p(i \mid s(i), t, \Hist_t)$
exactly, the log-likelihood of the marked stream splits into a ground term plus a mark term, and the
two factors can be estimated separately without loss.}

\begin{proof}
We prove the finite-mark-space case, which is the case used: the marks are the finitely many firms
of a sector. Represent the marked stream as the multivariate point process $(N_i)_{i \in
\mathcal{M}}$, $|\mathcal{M}| < \infty$, on the common internal history $(\Hist_t)$, each $N_i$
simple with predictable intensity $\lambda_i$ satisfying \eqref{eq:app-integrable}, and no two
components jumping together (simplicity of the marked stream).

\emph{Ground intensity.} $N_g = \sum_i N_i$ is a simple counting process and $N_g - \int_0^\cdot
\sum_i \lambda_i$ is a finite sum of martingales, hence a martingale; the integrand $\lambda_g :=
\sum_i \lambda_i$ is predictable and nonnegative, so $\lambda_g$ is the intensity of $N_g$.

\emph{The mark kernel.} Define $p(i \mid t) = \lambda_i(t)/\lambda_g(t)$ on $\{\lambda_g(t) > 0\}$
and let $p(\cdot \mid t)$ be an arbitrary fixed probability vector on $\{\lambda_g(t) = 0\}$. Then
$p(\cdot \mid t)$ is a probability vector for every $t$ and $\lambda_i = \lambda_g \, p(i \mid
\cdot)$ \emph{everywhere}: on $\{\lambda_g = 0\}$ both sides vanish because $0 \le \lambda_i \le
\lambda_g$. This is \eqref{eq:hw-factor}, an identity, with the firm-level intensity $\lambda_i$ and
$\lambda^g_s = \sum_{j : s(j) = s} \lambda_j$. Its interpretation as the conditional mark law is the
Campbell-type identity: for every bounded predictable $H(t, i)$, Lemma~\ref{lem:app-mart}(ii)
applied componentwise gives
\[
\E \sum_k H(t_k, m_k) \;=\; \sum_{i} \E \int_0^T H(t, i)\, \lambda_i(t)\, \dd t
\;=\; \E \int_0^T \lambda_g(t) \sum_i H(t, i)\, p(i \mid t)\, \dd t :
\]
conditionally on an event at $t$ and on $\Hist_{t^-}$, the mark is distributed as $p(\cdot \mid t)$
\cite[Ch.~6.4]{daley2003}.

\emph{Likelihood factorization.} For a realization $(t_1, m_1), \dots, (t_n, m_n)$ on $[0,T]$, the
likelihood of a multivariate point process with internal history is
\[
L_T \;=\; \Big[\prod_{k=1}^n \lambda_{m_k}(t_k)\Big]\,
\exp\Big(-\int_0^T \lambda_g(u)\, \dd u\Big),
\]
the product of successive conditional densities of (next time, next mark) given the past; we take
this representation from \cite[Ch.~7.2--7.3]{daley2003}, where it follows from the hazard
representation of the internal-history intensity. Substituting $\lambda_{m_k} = \lambda_g\, p(m_k
\mid \cdot)$,
\[
\log L_T \;=\; \underbrace{\sum_k \log \lambda_g(t_k) - \int_0^T \lambda_g\,\dd u}_{\text{ground
term}} \;+\; \underbrace{\sum_k \log p(m_k \mid t_k)}_{\text{mark term}} .
\]
If $\lambda_g$ carries parameters $\theta_1$ and $p$ carries parameters $\theta_2$ with $(\theta_1,
\theta_2)$ variation independent, the two terms share no parameters and joint maximization separates
into the two stage-wise maximizations: the split is lossless at population level, exactly as
Remark~\ref{rem:hw-lossless} asserts. For general (non-finite) mark spaces the same theorem holds
with $p$ a probability kernel; see \cite[Ch.~6--7]{daley2003}.
\end{proof}

\paragraph{Claim (Proposition~\ref{prop:hw-sampled}).}
\emph{Fix $m$ and replace, for each event, the risk set in \eqref{eq:hw-logit} by the winner plus a
uniform random subsample of $m$ non-winners drawn independently of utilities. The score of the
resulting conditional likelihood has mean zero at the true $w^*$, and its maximizer is consistent.}

\begin{proof}
Assume events are i.i.d.\ draws of (risk set $\RS$ with $|\RS| = r$, features $\{z_j\}_{j \in \RS}$,
winner $c$), with $\Prob(c = i \mid \RS, z) = e^{u_i}/\sum_{j \in \RS} e^{u_j}$, $u_j = w^{*\top}
z_j$; integrability $\E[\max_{j \in \RS} \|z_j\|] < \infty$; and the identification condition that
for every $v \ne 0$, the within-sample-set utilities $v^\top z_j$, $j \in D$, are non-constant with
positive probability. (In the sequential setting of Chapter~\ref{ch:hawkes} the same argument runs
conditionally on the event filtration, the score becomes a martingale difference array, and
consistency holds under ergodicity; we give the i.i.d.\ proof.)

\emph{Step 1: exact conditional likelihood (uniform conditioning).} Let $D = \{c\} \cup S$ where $S$
is a uniform $m$-subset of $\RS \setminus \{c\}$ drawn independently of everything else given $(\RS,
c)$. For a fixed $(m{+}1)$-subset $d \subseteq \RS$ and $i \in d$,
\[
\Prob(c = i,\ D = d \mid \RS, z)
= \frac{e^{u_i}}{\sum_{j \in \RS} e^{u_j}} \cdot \binom{r-1}{m}^{-1},
\]
and the second factor is the same for every $i \in d$ --- McFadden's uniform-conditioning property
\cite{mcfadden1974}. Dividing by the sum over $i' \in d$, $\Prob(c = i \mid D = d, \RS, z) = e^{u_i}
/ \sum_{j \in d} e^{u_j}$: the sampled softmax is an \emph{exact} conditional likelihood, not an
approximation.

\emph{Step 2: zero-mean score.} Let $\ell(w; c, D) = w^\top z_c - \log \sum_{j \in D} e^{w^\top
z_j}$, with gradient $z_c - \sum_j p_w(j \mid D) z_j$, bounded by $2\max_{j}\|z_j\|$, hence
integrable. For each fixed $d$, $\sum_{i \in d} p_{w^*}(i \mid d)\, \nabla_w \log p_w(i \mid
d)\,|_{w = w^*} = \nabla_w \sum_i p_w(i \mid d)\,|_{w^*} = 0$ (a finite sum, differentiated
termwise), so the score has conditional mean zero given $(D, \RS, z)$ at $w^*$, and therefore
unconditional mean zero.

\emph{Step 3: concavity and identification.} $\ell(\cdot\,; c, d)$ is concave: its Hessian is
$-\mathrm{Cov}_{p_w(\cdot \mid d)}(z) \preceq 0$ (log-sum-exp of affine functions is convex). Let
$L(w) = \E[\ell(w)]$. By Step 1 the model is correctly specified conditionally on $D$, so $L(w^*) -
L(w) = \E[\,\mathrm{KL}(p_{w^*}(\cdot \mid D)\, \|\, p_w(\cdot \mid D))\,] \ge 0$ (Gibbs' inequality
conditionally, then take expectations), with equality iff $p_w(\cdot \mid D) = p_{w^*}(\cdot \mid
D)$ a.s., i.e.\ iff $(w - w^*)^\top (z_i - z_j) = 0$ for all $i, j \in D$ a.s.\ --- excluded for $w
\ne w^*$ by the identification condition. So $w^*$ is the unique maximizer of $L$.

\emph{Step 4: consistency.} Let $\ell_n$ be the sample average criterion. For each fixed $w$,
$\ell_n(w) \to L(w)$ a.s.\ (strong law, integrability from Step 2). Pointwise almost-sure
convergence of concave functions on a countable dense subset of $\mathbb{R}^d$ implies uniform
convergence on compact sets (a standard fact of convex analysis). Fix $\epsilon > 0$ and let $\delta
= L(w^*) - \max_{\|w - w^*\| = \epsilon} L(w) > 0$ (positive by uniqueness and continuity).
Eventually, uniformly on the sphere, $\ell_n(w) \le L(w) + \delta/3 \le L(w^*) - 2\delta/3 \le
\ell_n(w^*) - \delta/3$; a concave function that is larger at the center than everywhere on a sphere
attains its maximum inside the ball, so every maximizer $\hat w_n$ satisfies $\|\hat w_n - w^*\| \le
\epsilon$ eventually, almost surely.
\end{proof}

\begin{lemma}[Gaussian moment generating function]
\label{lem:app-mgf}
If $z \sim \mathcal{N}(\mu, \Sigma)$ in $\mathbb{R}^q$ and $w \in \mathbb{R}^q$, then
$\E[\exp(w^\top z)] = \exp(w^\top \mu + \tfrac12 w^\top \Sigma w)$.
\end{lemma}

\begin{proof}
Factor $\Sigma = L L^\top$ (possible for any positive semidefinite matrix) and write $z
\stackrel{d}{=} \mu + L\varepsilon$ with $\varepsilon \sim \mathcal{N}(0, I_q)$. Then $w^\top z =
w^\top \mu + a^\top \varepsilon$ with $a = L^\top w$, and by independence of the coordinates $\E
e^{a^\top \varepsilon} = \prod_i \E e^{a_i \varepsilon_i}$. The univariate computation completes the
square: $\E e^{a_i \varepsilon_i} = (2\pi)^{-1/2} \int e^{a_i x - x^2/2}\, \dd x = e^{a_i^2/2}
(2\pi)^{-1/2} \int e^{-(x - a_i)^2/2}\, \dd x = e^{a_i^2/2}$. Multiply and use $\|a\|^2 = w^\top
\Sigma w$.
\end{proof}

\paragraph{Claim (Proposition~\ref{prop:hw-entrant}).}
\emph{With $N_{\mathrm{pool}}$ latent entrants whose features are i.i.d.\ $\mathcal{N}(\mu, \Sigma)$
independent of $\Hist_t$, replacing the entrant block of the softmax denominator by its expectation
yields the single pseudo-candidate score \eqref{eq:hw-entrant}.}

\begin{proof}
By linearity and Lemma~\ref{lem:app-mgf}, $\E \sum_{j=1}^{N_{\mathrm{pool}}} e^{w^\top z_j} =
N_{\mathrm{pool}}\, e^{w^\top \mu + w^\top \Sigma w/2}$, and matching $e^{q_{\mathrm{new}}}$ to this
expectation gives $q_{\mathrm{new}} = \log N_{\mathrm{pool}} + w^\top \mu + \tfrac12 w^\top \Sigma
w$, which is \eqref{eq:hw-entrant}. The accuracy of the mean-field replacement is quantified as
follows. Condition on $\Hist_t$ and let $A > 0$ be the (now fixed) funded-pool denominator $\sum_{j
\in \RS_{s,t}} e^{q_{j,t}}$ and $T = \sum_j e^{w^\top z_j}$ the entrant sum, independent of
$\Hist_t$ by assumption. Any funded candidate's choice probability is $e^{q_i}\, \E[(A + T)^{-1}]$.
The map $f(x) = (A + x)^{-1}$ on $[0, \infty)$ has $0 < f''(x) = 2 (A+x)^{-3} \le 2 A^{-3}$; a
second-order Taylor expansion of $f$ about $\E T$ with Lagrange remainder, followed by taking
expectations (the linear term has mean zero), gives
\[
\big|\, \E[f(T)] - f(\E T) \,\big| \;\le\; \frac{\Var(T)}{A^3},
\qquad
\Var(T) = N_{\mathrm{pool}}\; e^{2 w^\top\mu + w^\top\Sigma w}\big(e^{w^\top \Sigma w} - 1\big),
\]
the variance being the standard lognormal computation $\Var(e^X) = e^{2m + v}(e^v - 1)$ for $X \sim
\mathcal{N}(m, v)$, from Lemma~\ref{lem:app-mgf} applied at $w$ and $2w$. When $A$ and $\E T$ both
grow like the pool size, $f \asymp N^{-1}$ while the error is $O(N^{-2})$: the relative error of the
mean-field step is $O(1/N_{\mathrm{pool}})$, negligible at the pool sizes of
Chapter~\ref{ch:hawkes}.
\end{proof}

\begin{proposition}[Pool-size non-identifiability, made precise]
\label{prop:app-poolident}
In any parametrization in which the entrant block enters the choice probabilities only through the
scalar score $q_{\mathrm{new}} = b_0 + \log N_{\mathrm{pool}} + w^\top\mu + \tfrac12 w^\top \Sigma
w$ with a constant pool, the likelihood of every dataset is invariant under $(b_0, \log
N_{\mathrm{pool}}) \mapsto (b_0 + c, \log N_{\mathrm{pool}} - c)$ for every real $c$: the constant
$b_0$ and the pool size (equivalently, the pool size and the entrant mean utility) are not
separately identified. If instead $N_{\mathrm{pool}, t}$ varies over events and enters as a known
offset, $q_{\mathrm{new}, t} = b_0 + \log N_{\mathrm{pool}, t} + \cdots$, then $b_0$ is identified.
\end{proposition}

\begin{proof}
Every choice probability --- of a funded candidate or of the newcomer alternative --- depends on the
entrant parameters only through $e^{q_{\mathrm{new}}}$ in the softmax; the stated transformation
leaves $q_{\mathrm{new}}$, hence every probability, hence the likelihood of any dataset, unchanged.
This is exact non-identification, not a flat-likelihood approximation. With known offsets, two
parameter values produce identical probabilities only if $b_0 + \log N_{\mathrm{pool},t} = b_0' +
\log N_{\mathrm{pool},t}$ for all $t$, forcing $b_0 = b_0'$. This is why the production checklist of
Chapter~\ref{ch:selection} requires recalibrating the newcomer constant whenever the pool definition
changes: the fitted constant is a statement about $b_0 + \log N_{\mathrm{pool}}$, never about
entrant quality alone.
\end{proof}

\section{Proofs for Chapter~\ref{ch:censored}}
\label{app:proofs-ch9}

\paragraph{Claim (Proposition~\ref{prop:cen-volterra}).}
\emph{For the linear Hawkes process \eqref{eq:hw-intensity} with locally bounded baseline $s \ge 0$
and integrable kernel $\phi \ge 0$, the mean intensity $\xi(t) = \E[\lambda(t)]$ is locally
integrable and solves the Volterra equation \eqref{eq:cen-volterra}; it is its unique locally
integrable solution. The multivariate statement holds componentwise.}

\begin{proof}
Realize the process by the cluster construction in the proof of Theorem~\ref{thm:hw-cluster}
(subcritical case) or note that \eqref{eq:app-integrable} holds by
Lemma~\ref{lem:app-branching}(iii) for every $\kappa$; either way $\E N(T) < \infty$ for all $T$, so
Lemma~\ref{lem:app-fubini} applies. Writing the history term as an integral against $\dd N$ and
applying Lemma~\ref{lem:app-fubini} with the nonnegative measurable $f(u) = \phi(t-u)\ind{u < t}$,
\[
\xi(t) = s(t) + \E \int_{(0,t)} \phi(t-u)\, \dd N(u)
= s(t) + \E \int_0^t \phi(t-u)\, \lambda(u)\, \dd u
= s(t) + \int_0^t \phi(t-u)\, \xi(u)\, \dd u ,
\]
the last step by Tonelli (everything nonnegative). Local integrability: $\int_0^T \xi = \E
\Lambda(T) = \E N(T) < \infty$, again by Lemma~\ref{lem:app-fubini} with $f \equiv 1$; in particular
$\xi$ is finite a.e.\ and the equation holds with finite right-hand side a.e. Uniqueness: if $\xi_1,
\xi_2$ are locally integrable solutions, $d = |\xi_1 - \xi_2|$ satisfies $d \le \phi * d$ a.e.,
hence by iteration $d \le \phi^{*n} * d$ for all $n$; integrating over $[0,T]$ and using $\int_0^T
(\phi^{*n} * d) \le (\int_0^T \phi^{*n})(\int_0^T d)$ with $\int_0^T \phi^{*n} = \kappa^n \Prob(S_n
\le T) \to 0$ (Lemma~\ref{lem:app-branching}(iii)) forces $\int_0^T d = 0$ for every $T$. The
linearity of \eqref{eq:hw-intensity} is what lets the expectation pass through --- for the log-link
models of Chapter~\ref{ch:hawkes} the first step fails, exactly as stated there. For the
multivariate form see also \cite{bacry2015}.
\end{proof}

\paragraph{Claim (Proposition~\ref{prop:cen-rational}).}
\emph{The MBPP map $s \mapsto \xi$ of \eqref{eq:cen-volterra} is realizable as a finite-dimensional
time-invariant linear ODE with output $\xi$ iff the kernel's Laplace transform $\hat\phi(z)$ is a
strictly proper rational function, equivalently iff $\phi$ is a finite sum of
polynomial-times-exponential terms (complex-conjugate pairs allowed).}

\begin{proof}
Call the map \emph{realizable} if there exist finite $n$ and $(A, b, c)$, $A \in \mathbb{R}^{n
\times n}$, $b, c \in \mathbb{R}^n$, such that for every locally bounded input $s$, the solution of
$x' = A x + b\, \xi$, $x(0) = 0$, with $\xi = s + c^\top x$, reproduces the Volterra solution.
Substituting the output equation, this is the ordinary ODE $x' = (A + b c^\top) x + b s$ with output
$\xi = c^\top x + s$: a finite linear system driven by $s$.

\emph{Realizability is a property of the kernel.} By variation of constants, the feedback loop
computes $c^\top x(t) = \int_0^t g(t-u)\, \xi(u)\, \dd u$ with $g(t) = c^\top e^{At} b$
(differentiate to verify; $e^{At}$ is the matrix exponential). The system solves
\eqref{eq:cen-volterra} for all inputs iff $\int_0^t (g - \phi)(t - u)\, \xi(u)\, \dd u = 0$ for all
$t$ and all solutions $\xi$; since any continuous $\xi$ arises from the input $s = \xi - \phi *
\xi$, the du Bois-Reymond argument (test against approximate identities) forces $\phi = g$ a.e. So
realizability means precisely $\phi(t) = c^\top e^{At} b$ a.e.\ for some finite triple.

\emph{Realizable kernels have strictly proper rational transforms.} For $\mathrm{Re}\, z > \|A\|$
the integral $\int_0^\infty e^{-zt} e^{At}\, \dd t$ converges absolutely and equals $(zI - A)^{-1}$
(multiply by $zI - A$ and integrate by parts, or integrate the exponential series termwise). Hence
$\hat\phi(z) = c^\top (zI - A)^{-1} b = c^\top \mathrm{adj}(zI - A)\, b / \det(zI - A)$: a ratio of
a polynomial of degree at most $n - 1$ to one of degree $n$ --- rational and strictly proper.

\emph{Conversely.} Let $\hat\phi = q/p$ be strictly proper rational. Partial fractions over
$\mathbb{C}$: $\hat\phi(z) = \sum_j \sum_{r=1}^{r_j} \alpha_{jr} (z - z_j)^{-r}$. Termwise, $(z -
z_j)^{-r}$ is the Laplace transform of $t^{r-1} e^{z_j t}/(r-1)!$ (direct integration for
$\mathrm{Re}\, z > \mathrm{Re}\, z_j$), so the polynomial-times-exponential function $g(t) =
\sum_{j,r} \alpha_{jr}\, t^{r-1} e^{z_j t}/(r-1)!$ has $\hat g = \hat\phi$; by injectivity of the
Laplace transform on locally integrable functions of exponential order (a classical uniqueness
theorem), $\phi = g$ a.e. Realization: for a single term, the $r \times r$ Jordan block $J$ at $z_j$
satisfies $e^{Jt} = e^{z_j t}\,(t^{k-l}/(k-l)!)_{l \le k}$ (upper triangular), so $e_1^\top e^{Jt}
e_r = t^{r-1} e^{z_j t}/(r-1)!$; a direct sum of such blocks (conjugate pairs combined into real $2r
\times 2r$ blocks) with the matching weights in $c$ gives $\phi(t) = c^\top e^{At} b$. Finally, the
exponential kernel recovers the scalar ODE \eqref{eq:cen-ode}: $\phi = \kappa\theta e^{-\theta t}$
is realized by $A = -\theta$, $b = 1$, $c = \kappa\theta$, and the closed loop $x' = (-\theta +
\kappa\theta) x + s\kappa\theta/( \kappa\theta)$ --- written on $h = c x$ --- is $h' =
\kappa\theta\, s - \theta(1-\kappa)\, h$, exactly \eqref{eq:cen-ode}. Power-law kernels have
branch-cut (non-rational) transforms and admit no finite realization, which is why they require
quadrature or exponential-sum approximation, as stated in Section~\ref{sec:cen-solve}.
\end{proof}

We now make the $\kappa$--$\theta$ ridge quantitative. Recall the closed-form bucket objects of
Section~\ref{sec:cen-worked}: with equal buckets of width $\Delta$, piecewise-constant baseline $s_w
\in [s_{\min}, s_{\max}]$, $0 < s_{\min}$, relaxation rate $\gamma = \theta(1 - \kappa)$ and $\rho =
e^{-\gamma\Delta}$, the state $h_w = h(o_{w-1})$ obeys \eqref{eq:cen-state} with target levels
$h_w^\ast = \kappa s_w/(1-\kappa)$, and the bucket compensator is \eqref{eq:cen-xiw}. (Both displays
follow from \eqref{eq:cen-ode} by solving the linear ODE on a bucket where $s$ is constant and
integrating $\xi = s + h$; \eqref{eq:cen-ode} itself follows from \eqref{eq:cen-volterra} by writing
$h(t) = \kappa\theta e^{-\theta t}\int_0^t e^{\theta u} \xi(u)\, \dd u$ and differentiating.)

\begin{proposition}[The limiting singularity of the censored information matrix]
\label{prop:app-ridge}
Fix $\kappa \in (0,1)$ and the baseline range $0 < s_{\min} \le s_w \le s_{\max}$, take the
stationary initialization $h_1 = h_1^\ast$, and let $C_0 = \kappa(s_{\max} - s_{\min})/(1-\kappa)$.
Then for all $w$: (i) $|h_w - h_w^\ast| \le C_0$ and $\Xi_w = \Delta s_w/(1-\kappa) + R_w$ with
$|R_w| \le C_0/\gamma$; (ii) $\big|\partial \Xi_w / \partial\theta\big| \le C_0 (1-\kappa)\, \big(
2\Delta\rho/\gamma + \gamma^{-2} \big)$; (iii) $\partial \Xi_w / \partial\kappa = \Delta s_w
(1-\kappa)^{-2} + r_w$ with $|r_w| \le C_1 (\Delta \rho + 1/\gamma)$, where $C_1$ depends only on
$(\kappa, s_{\min}, s_{\max})$. Consequently, whenever $\gamma \ge \gamma_0 > 0$ and $\gamma\Delta
\to \infty$,
\[
\frac{\sup_w |\partial \Xi_w/\partial\theta|}{\inf_w \partial \Xi_w/\partial\kappa}
\;\longrightarrow\; 0,
\]
and the per-bucket Fisher information of the Poisson likelihood \eqref{eq:cen-lik} in $(\kappa,
\theta)$, $I = \sum_w \Xi_w^{-1} (\nabla \Xi_w)(\nabla \Xi_w)^\top$, becomes singular:
$I_{\theta\theta}$ and $I_{\kappa\theta}$ vanish relative to $I_{\kappa\kappa}$, and the profile
information for $\theta$, $I_{\theta\theta} - I_{\theta\kappa}^2/I_{\kappa\kappa}$, tends to zero
per bucket. In the regime $\theta \to \infty$ with $\Delta$ fixed, $\partial \Xi_w/\partial\theta
\to 0$ absolutely while $\partial \Xi_w/\partial\kappa \to \Delta s_w (1-\kappa)^{-2}$, bounded away
from zero: $\theta$ leaves the first-moment likelihood entirely.
\end{proposition}

\begin{proof}
\emph{Fisher information formula.} For independent $C_w \sim \mathrm{Poisson}(\Xi_w(\eta))$ the
score in a parameter $\eta$ is $\sum_w (C_w/\Xi_w - 1)\, \partial_\eta \Xi_w$; the summands are
independent with mean zero and variance $(\partial_\eta \Xi_w)^2/\Xi_w$ (since $\Var C_w = \Xi_w$),
giving the stated $I$.

(i) The levels $h_w^\ast$ lie in $[\kappa s_{\min}/(1-\kappa),\, \kappa s_{\max}/(1-\kappa)]$, an
interval of length $C_0$. By \eqref{eq:cen-state}, $h_{w+1}$ is a convex combination of $h_w$ and
$h_w^\ast$; with $h_1 = h_1^\ast$ in the interval, induction keeps every $h_w$ in it, so $|h_w -
h_w^\ast| \le C_0$. Then $R_w = (h_w - h_w^\ast)(1-\rho)/\gamma$ obeys $|R_w| \le C_0/\gamma$.

(ii) Only $\gamma$ and $\rho$ carry $\theta$; $h_w^\ast$ does not. Differentiating
\eqref{eq:cen-state} in $\theta$, the derivative state $D_w = \partial h_w/\partial\theta$ satisfies
$D_1 = 0$ and $D_{w+1} = \rho D_w + (h_w - h_w^\ast)\, \partial\rho/\partial\theta$ with
$\partial\rho/\partial\theta = -(1-\kappa)\Delta\rho$, so $|D_{w+1}| \le \rho |D_w| +
(1-\kappa)\Delta\rho\, C_0$ and by the geometric recursion $|D_w| \le (1-\kappa)\Delta\rho\,
C_0/(1-\rho)$ for all $w$. Also $\partial_\theta [(1-\rho)/\gamma] =
(1-\kappa)\big[\Delta\rho/\gamma - (1-\rho)/\gamma^2\big]$, bounded in modulus by
$(1-\kappa)(\Delta\rho/\gamma + \gamma^{-2})$. Hence
\[
\Big|\frac{\partial \Xi_w}{\partial\theta}\Big|
\le |D_w|\, \frac{1-\rho}{\gamma} + C_0 (1-\kappa)\Big(\frac{\Delta\rho}{\gamma} +
\frac{1}{\gamma^2}\Big)
\le C_0 (1-\kappa) \Big( \frac{2\Delta\rho}{\gamma} + \frac{1}{\gamma^2} \Big).
\]
(iii) The leading term of $\Xi_w$ differentiates to $\Delta s_w (1-\kappa)^{-2}$. For the remainder,
$\partial_\kappa h_w^\ast = s_w/(1-\kappa)^2 \le s_{\max}(1-\kappa)^{-2}$, and the derivative state
$F_w = \partial h_w/\partial\kappa$ obeys $F_{w+1} = (1-\rho)\, \partial_\kappa h_w^\ast + \rho F_w
+ (h_w - h_w^\ast)\, \partial\rho/\partial\kappa$ with $\partial\rho/\partial\kappa =
\theta\Delta\rho$, whence, using $\theta = \gamma/(1-\kappa)$ and the geometric recursion again,
$|F_w| \le s_{\max}(1-\kappa)^{-2} + C_0\,\Delta\rho\,\gamma\,[(1-\kappa)(1-\rho)]^{-1}$. Combining
with $\partial_\kappa[(1-\rho)/\gamma] = \theta\big[\Delta\rho/\gamma - (1-\rho)/\gamma^2\big]$,
every term of $\partial_\kappa R_w$ is bounded by a constant multiple (depending only on $\kappa,
s_{\min}, s_{\max}$) of $\Delta\rho$ or $1/\gamma$, which is (iii).

The displayed ratio limit follows by dividing (ii) by (iii): under $\gamma \ge \gamma_0$ and
$\gamma\Delta \to \infty$, the numerator terms $\Delta\rho/\gamma$ and $\gamma^{-2}$ are both
$o(\Delta)$ ($\Delta\rho \to 0$ superexponentially, and $\gamma^{-2} \le (\gamma_0
\gamma\Delta)^{-1}\gamma\Delta\cdot\gamma^{-1}\ldots$ more simply $\gamma^{-2}/\Delta =
(\gamma\Delta)^{-1}\gamma^{-1} \le (\gamma\Delta)^{-1}\gamma_0^{-1} \to 0$), while the denominator
is bounded below by $\Delta s_{\min}(1-\kappa)^{-2}$ minus a vanishing correction. Since $\Xi_w$
stays in a fixed positive interval around $\Delta s_w/(1-\kappa)$ by (i), the information entries
inherit the same ratios: $I_{\theta\theta}/I_{\kappa\kappa} \to 0$,
$|I_{\kappa\theta}|/I_{\kappa\kappa} \to 0$, and the profile information for $\theta$ is at most
$I_{\theta\theta} \to 0$ per bucket. By the Cramer--Rao bound, the variance of any unbiased
estimator of $\theta$ from bucket means diverges in this limit. The $\theta \to \infty$, $\Delta$
fixed statement reads off (ii) and (iii) directly.
\end{proof}

\paragraph{Claim (Proposition~\ref{prop:cen-ridge}).}
\emph{In the stationary regime with slowly varying baseline and bucket width $\Delta$: (i) $\partial
\Xi_w/\partial\kappa$ is an order-one level effect of size $\Xi_w/(1-\kappa)$; (ii) $\theta$ enters
$\Xi_w$ only through transient terms of relative size governed by $e^{-\theta(1-\kappa)\Delta}$, so
the Fisher information of \eqref{eq:cen-lik} is nearly singular along a ridge in the $(\kappa,
\theta)$ plane, worsening monotonically with $\Delta$; as $\theta(1-\kappa)\Delta \to \infty$,
$\theta$ leaves the first-moment likelihood entirely.}

\begin{proof}
This is Proposition~\ref{prop:app-ridge} restated qualitatively. For (i): $\Delta s_w
(1-\kappa)^{-2} = [\Delta s_w/(1-\kappa)] \cdot (1-\kappa)^{-1} = \Xi_w/(1-\kappa)$ up to the
$O(1/\gamma)$ correction of part (i) of that proposition. For (ii): the bounds in part (ii) are
proportional to $\Delta\rho/\gamma$ and $\gamma^{-2}$ with $\rho = e^{-\theta(1-\kappa)\Delta}$,
both vanishing relative to the $\kappa$-score as $\gamma\Delta \to \infty$; the sign oscillation
across buckets follows from the derivative state recursion for $D_w$, which alternates in sign as
the lagged response $h_w - h_w^\ast$ changes sign with the baseline steps. The near-singularity of
the information matrix and the escape of $\theta$ from the likelihood are the two displayed
conclusions of Proposition~\ref{prop:app-ridge}. What censored data retain of $\theta$ beyond the
first moment lives in autocovariances, which the MBPP likelihood discards by construction ---
Remark~\ref{rem:cen-banana} and the reporting rules of Table~\ref{tab:cen-worked} are the
operational consequences.
\end{proof}

\chapter{Schema reference}
\label{app:schema}

This appendix is the lookup companion to Chapter~\ref{ch:data}: field-level types and typical missingness
for the three PitchBook relations, a macro-feed sample showing the publish-lag structure, and coverage by
deal class. Missingness figures are representative magnitudes for a US-heavy export, matching the
calibration of the synthetic observation layer; re-profile any given export before modelling.

\section{The deals table}\label{app:schema-deals}

\begin{table}[htbp]
\centering
\footnotesize
\begin{tabular}{llp{5.0cm}}
\toprule
Field & Type & Typical missingness / notes \\
\midrule
\multicolumn{3}{l}{\emph{Identifiers}} \\
\texttt{company\_id}, \texttt{deal\_id} & string key & none; join keys \\
\texttt{company\_name} & string & none; display only, never a key \\
\texttt{deal\_date} & date & none, but jitter up to 6 days \\
\midrule
\multicolumn{3}{l}{\emph{Industry and universe}} \\
\texttt{primary\_industry\_sector} & categorical & $\sim$1\%; $\sim$3\% mislabelled \\
\texttt{primary\_industry\_group}, \texttt{primary\_industry\_code} & categorical & $\sim$2\% \\
\texttt{company\_universe} & categorical & none; snapshot semantics \\
\texttt{current\_financing\_status} & categorical & snapshot --- forbidden as feature \\
\midrule
\multicolumn{3}{l}{\emph{Deal terms}} \\
\texttt{deal\_type}, \texttt{deal\_class}, \texttt{deal\_status} & categorical & none \\
\texttt{deal\_type2}, \texttt{deal\_type3} & categorical & 40\%, 70\% \\
\texttt{deal\_size} & float (MUSD) & 28\% undisclosed; 15\% of rest estimated (18\% rel.\ error) \\
\texttt{deal\_size\_status} & categorical & none; flags estimated sizes \\
\texttt{pre\_money\_valuation}, \texttt{post\_valuation} & float (MUSD) & 55\% \\
\texttt{post\_valuation\_status} & categorical & follows valuations \\
\texttt{pct\_acquired} & float & 60\%; mostly M\&A rows \\
\texttt{vc\_round}, \texttt{vc\_round\_up\_down\_flat} & categorical & VC only; 65\% for up/down/flat \\
\texttt{type\_of\_stock}, \texttt{sellers} & string & 70\%+ \\
\midrule
\multicolumn{3}{l}{\emph{Capital structure}} \\
\texttt{total\_invested\_capital}, \texttt{total\_invested\_equity} & float & 45\% \\
\texttt{debt}, \texttt{total\_new\_debt}, \texttt{total\_debt} & float & 80\% on VC; better on debt \\
\texttt{raised\_to\_date}, \texttt{total\_raised\_to\_date} & float & 30\%; recompute from history \\
\texttt{total\_investor\_ownership} & float & 75\% \\
\midrule
\multicolumn{3}{l}{\emph{Financials at deal}} \\
\texttt{revenue}, \texttt{ebitda} & float & 70\%, 80\% (US); far lower in EU \\
\texttt{gross\_profit}, \texttt{net\_income} & float & 85\% \\
\texttt{valuation\_ebitda}, \texttt{valuation\_revenue}, \texttt{valuation\_cash\_flow} & float & 85\%+ \\
\texttt{deal\_size\_ebitda}, \texttt{deal\_size\_revenue}, \texttt{deal\_size\_cash\_flow} & float & 85\%+ \\
\texttt{revenue\_growth\_since\_last\_deal} & float & 80\% \\
\texttt{number\_of\_employees} & int & 40\% \\
\midrule
\multicolumn{3}{l}{\emph{Firmography carried on deal rows}} \\
\texttt{year\_founded} & int & 5\% \\
\texttt{hq\_location}, \texttt{hq\_country}, \texttt{hq\_region}, \texttt{hq\_sub\_region} & string & 2\% \\
\bottomrule
\end{tabular}
\caption{Deals table: fields, types, and representative missingness.}
\label{tab:app-deals}
\end{table}

\section{The companies and investors tables}\label{app:schema-companies}

\begin{table}[htbp]
\centering
\footnotesize
\begin{tabular}{llp{5.4cm}}
\toprule
Field & Type & Notes \\
\midrule
\texttt{company\_id} & string key & join key; 12\% of real firms untracked \\
\texttt{description} & string & free text; 8\% missing \\
\texttt{year\_founded} & int & 5\% missing \\
\texttt{employees} & int & \emph{current} snapshot; forbidden as historical feature \\
\texttt{financing\_status}, \texttt{business\_status}, \texttt{ownership\_status} & categorical & snapshots \\
\texttt{total\_money\_raised} & float & snapshot; recompute from deals \\
\texttt{last\_financing\_date}, \texttt{last\_financing\_size}, \texttt{last\_financing\_valuation} & mixed & snapshots of latest deal \\
\midrule
\texttt{investor\_id} & string key & investors relation \\
\texttt{investorname} & string & display only \\
\texttt{investorstatus} & categorical & current snapshot \\
\texttt{investorsince} & date & dated; usable point-in-time \\
\bottomrule
\end{tabular}
\caption{Companies table (top block) and investors relation (bottom block). Snapshot fields may validate
history but must never enter features (Chapter~\ref{ch:data}).}
\label{tab:app-companies}
\end{table}

\section{Macro feed sample}\label{app:schema-macro}

\begin{table}[htbp]
\centering
\small
\begin{tabular}{llll}
\toprule
\texttt{series\_code} & \texttt{ref\_date} & \texttt{publish\_date} & \texttt{value} \\
\midrule
\texttt{US\_GDP\_QOQ\_SAAR} & 2023-12-31 & 2024-01-25 & 3.3 \\
\texttt{US\_CPI\_YOY} & 2024-01-31 & 2024-02-13 & 3.1 \\
\texttt{US\_UNRATE} & 2024-01-31 & 2024-02-02 & 3.7 \\
\texttt{US\_10Y\_YIELD} & 2024-02-15 & 2024-02-15 & 4.24 \\
\texttt{EA\_GDP\_QOQ} & 2023-12-31 & 2024-01-30 & 0.0 \\
\texttt{US\_PE\_DRYPOWDER} & 2023-12-31 & 2024-03-15 & 955.0 \\
\bottomrule
\end{tabular}
\caption{Macro feed sample rows. The \texttt{ref\_date}-to-\texttt{publish\_date} gap --- zero for market
series, weeks to months for releases --- is the point-in-time problem of Chapter~\ref{ch:us}: all joins
key on \texttt{publish\_date}.}
\label{tab:app-macro}
\end{table}

\section{Field coverage by deal class}\label{app:schema-coverage}

\begin{table}[htbp]
\centering
\small
\begin{tabular}{lccc}
\toprule
Field group & VC & M\&A & Debt \\
\midrule
\texttt{deal\_size} disclosed & 65\% & 45\% & 75\% \\
Valuations (\texttt{pre\_money}, \texttt{post}) & 55\% & 20\% & 5\% \\
\texttt{pct\_acquired} & 15\% & 70\% & 10\% \\
\texttt{vc\_round} fields & 95\% & --- & --- \\
Debt / capital-structure fields & 15\% & 35\% & 90\% \\
Financials at deal & 20\% & 45\% & 40\% \\
\bottomrule
\end{tabular}
\caption{Approximate non-missing coverage by deal class. No field is well populated across all classes;
class-conditional models or class indicators are mandatory (Chapter~\ref{ch:data}).}
\label{tab:app-coverage}
\end{table}

\chapter{Code map and quickstart}
\label{app:code}

\section{Directory layout}\label{app:code-map}

The accompanying code lives under \texttt{handbook/code}. The core package implements the estimators; the
\texttt{synthetic} package implements the two testbed worlds of Chapter~\ref{ch:selection} and their
observation-noise layer.

\begin{verbatim}
handbook/code/
  hawkes_calibration/        core package
    eventtime/               event-time Hawkes: kernels, MLE, thinning simulation
    mbpp/                    Volterra solver, interval-censored likelihood
    models/                  sector_ranker.py, sector_hazard.py,
                             sector_survival.py, sector_backtest.py
    operators/               pino.py (operator surrogate), training loops
  synthetic/
    macro.py                 latent macro factors and published series
    firms.py                 firm universe, entry and exit
    marks.py                 deal marks: size, valuation, stage
    regime_a.py              weekly two-stage DGP (well specified)
    regime_b.py              daily firm-level block-Hawkes DGP (misspecified)
    observe.py               noise layer: missingness, jitter, duplicates
    loaders.py               load_dataset, build_choice_sets
    fast_fit.py              fit_sector_glm_fast, fit_choice_fast
    generate.py              command-line entry point
  EXPERIMENTS.md             experiment briefings, seeds, expected outputs
  REVIEW.md                  review notes and open issues
\end{verbatim}

\section{Quickstart}\label{app:code-quickstart}

Generate both synthetic worlds (seed 0 reproduces every number in Chapter~\ref{ch:selection}):

\begin{verbatim}
cd handbook/code
PYTHONPATH=. python -m synthetic.generate --regime A --seed 0
PYTHONPATH=. python -m synthetic.generate --regime B --seed 0
\end{verbatim}

Load, build choice sets, and fit both stages:

\begin{verbatim}
from synthetic.loaders import load_dataset, build_choice_sets
from synthetic.fast_fit import fit_sector_glm_fast, fit_choice_fast

ds  = load_dataset("A", seed=0)
glm = fit_sector_glm_fast(ds)          # Stage 1: sector count GLM
cs  = build_choice_sets(ds, n_candidates=64, cooldown_weeks=26)
cl  = fit_choice_fast(cs)              # Stage 2: conditional logit
\end{verbatim}

\section{Regenerating the numbers}\label{app:code-regen}

Every table in Chapter~\ref{ch:selection} regenerates from these entry points:
Tables~\ref{tab:sel-stage1} and~\ref{tab:sel-choice} from the two fits above on both regimes;
Table~\ref{tab:sel-riskset} by passing naive versus oracle membership masks to
\texttt{build\_choice\_sets}; Table~\ref{tab:sel-stability} from the simulation guard in
\texttt{hawkes\_calibration/models}. \texttt{EXPERIMENTS.md} lists each experiment with its command, seed,
runtime, and the assertion it must satisfy --- including the causal-validation hypotheses ---
and \texttt{REVIEW.md} records known deviations; if a regenerated number drifts, trust the code.
